% 根据已评审 Markdown 需求基线排版；需求语义与编号保持一致。
\clearpage
\section{具体需求}\label{sec:3-具体需求}

\Needspace{10\baselineskip}
\subsection{外部接口与交互要求}\label{sec:31-外部接口与交互要求}

\Needspace{12\baselineskip}
\subsubsection{用户界面}\label{sec:311-用户界面}

系统应通过浏览器提供以下逻辑页面和操作入口。页面应体现用户身份与数据可见范围；具体 URL 和视觉布局由界面设计确定，公开商品和文章详情必须具有可直接访问的独立地址。

\begingroup
\smalltable
\begin{longtable}{>{\raggedright\arraybackslash}p{\dimexpr 0.150\linewidth-2\tabcolsep\relax}>{\raggedright\arraybackslash}p{\dimexpr 0.430\linewidth-2\tabcolsep\relax}>{\raggedright\arraybackslash}p{\dimexpr 0.420\linewidth-2\tabcolsep\relax}}
\caption{逻辑页面与操作入口}\label{tab:11}\\
\toprule
{\bfseries 区域} & {\bfseries 页面或入口} & {\bfseries 数据来源} \\
\midrule
\endfirsthead
\multicolumn{3}{r}{\footnotesize 表~\thetable~（续）}\\
\toprule
{\bfseries 区域} & {\bfseries 页面或入口} & {\bfseries 数据来源} \\
\midrule
\endhead
\midrule
\multicolumn{3}{r}{\footnotesize 续下页}\\
\endfoot
\bottomrule
\endlastfoot
品牌官网 & 首页、品牌介绍 & 后台品牌设置、Banner、推荐商品 \\
\tabledivider
商品 & 列表、详情 & 商品、分类、价格、库存、资料 \\
\tabledivider
内容支持 & 玩法列表和详情、FAQ、下载中心 & 文章、FAQ、文件记录 \\
\tabledivider
渠道与联系 & 经销商列表、联系表单 & 公开渠道、咨询记录 \\
\tabledivider
身份 & 注册、登录 & 账户与会话 \\
\tabledivider
零售 & 购物车、结算、模拟付款 & 购物车、订单与库存 \\
\tabledivider
个人中心 & 资料、订单、我的咨询、经销申请记录 & 当前用户的业务数据 \\
\tabledivider
合作申请 & 经销商申请 & 企业申请和审核记录 \\
\tabledivider
经销商中心 & 专属目录、下载、询价列表和详情 & 有效经销商可见数据 \\
\tabledivider
管理后台 & 总览、商品、内容、渠道、订单、用户、申请、工单、操作记录 & 受管理员权限保护的业务数据 \\
\tabledivider
辅助页面 & 使用说明、隐私说明、无权限、未找到、服务异常 & 固定说明和错误状态 \\
\end{longtable}
\endgroup

主要表单应具有可见标签、输入约束和明确反馈；涉及取消、确认收货等改变业务状态的动作应提供确认步骤。加载中、空结果、无权限、资源不存在和服务异常应具有不同反馈，详见第 3.1.3 节与 NFR-03。

\Needspace{12\baselineskip}
\subsubsection{软件与通信接口}\label{sec:312-软件与通信接口}

\begin{enumerate}
\item 前台、个人中心、经销商中心和后台应通过受控业务接口读写数据。身份、角色、数据归属、价格、库存和状态校验由服务端独立执行。
\item 接口应支持结构化请求与响应。列表应按对应 FR 提供分页、排序或筛选；分页列表应返回总条数，单页上限为 50。
\item 受保护接口应识别当前会话，并在处理请求时检查账户和合作状态。未获授权的字段不得包含在响应中。
\item 文件接口应返回正确文件类型；非公开文件下载应进行当次授权检查。文件格式、大小与替换规则按 FR-07、FR-29 执行。
\item 订单、申请、工单和询价创建接口应接受可用于重复请求识别的业务标识；重试与冲突处理按 BR-06 执行。
\item 本版本不调用外部支付、物流、邮件、地图、ERP、WMS 或 CRM 服务。模拟付款和发货产生的记录不得被解释为外部服务已执行。
\item 系统无专用硬件接口要求；客户端仅需具备受支持的浏览器与网络连接。
\end{enumerate}

\Needspace{12\baselineskip}
\subsubsection{输入校验与错误反馈}\label{sec:313-输入校验与错误反馈}

\begin{enumerate}
\item 字段错误显示在对应输入项附近，包含错误原因与修正方式；保留已经填写的非敏感内容。
\item 提交时显示处理状态并暂时阻止重复点击；成功后给出业务编号、结果和后续入口；失败后允许修正与重试。
\item 列表和详情区分加载中、无数据、加载失败、无权限和不存在，不用同一个空白界面代替。
\item 返回明确业务错误码及用户可读信息；字段错误可返回字段名与提示；服务端异常只显示通用说明和排查编号，不暴露堆栈、SQL 或文件路径。
\item 接口应采用一致的 HTTP 语义：请求格式错误 400、未认证 401、无权限 403、公开资源不存在或私人资源不可访问 404、状态／库存冲突 409、字段校验失败 422、请求受限 429、系统异常 500。具体业务错误码由接口说明统一定义。
\item 内容输入按普通文本转义；采用富文本时，对允许的标签和属性进行清洗。对文件名、链接、查询参数同样执行服务端验证。
\end{enumerate}

\Needspace{10\baselineskip}
\subsection{功能需求}\label{sec:32-功能需求}

以下功能需求均为 P0。每项需求包含需求陈述、业务理由、适用条件、行为约束、验收条件和验证方法；关联 TC 的完整场景见附录 A。涉及金额、库存、状态和重复请求时，同时适用第 3.3 节规则。

\Needspace{12\baselineskip}
\subsubsection{公开网站}\label{sec:321-公开网站}

\reqheading{FR-01}{FR-01 全站导航与基础页面}

\begin{enumerate}
\item {\bfseries 需求陈述}：系统应向访问者提供一致的站点导航、身份入口和基础状态页面。
\item {\bfseries 业务理由}：让用户能够识别当前位置并到达相应业务入口。
\item {\bfseries 适用条件}：所有访问者；公开页面不要求登录。
\item {\bfseries 行为与约束}：
\begin{enumerate}
\item {\bfseries 功能}：提供首页、产品、玩法与支持、经销商、关于、联系入口；显示当前登录状态，登录后可进入相应个人中心；提供购物车入口。
\item {\bfseries 规则}：移动端导航可展开与收起；当前页面有清楚的标题；商品和文章详情可返回列表；页脚提供使用说明及隐私说明。
\item {\bfseries 异常}：不存在的公开资源进入未找到页面；无权限资源给出登录或无权限提示；页面不得因接口失败永久停留在加载动画。
\end{enumerate}

\begin{samepage}
\item {\bfseries 验收条件}：桌面和手机均可从首页到达各个 P0 入口，链接无错误跳转；直接访问详情 URL 可打开对应页面。
\item {\bfseries 验证方法}：接口测试与业务场景验证；关联 TC-01、TC-38。
\end{samepage}
\end{enumerate}

\reqheading{FR-02}{FR-02 首页与品牌展示}

\begin{enumerate}
\item {\bfseries 需求陈述}：系统应展示由后台维护的品牌信息、首页内容和推荐商品。
\item {\bfseries 业务理由}：使品牌与商品信息能够随运营更新保持一致。
\item {\bfseries 适用条件}：所有访问者；内容由管理员维护并发布。
\item {\bfseries 行为与约束}：
\begin{enumerate}
\item {\bfseries 功能}：展示品牌定位、主 Banner、产品分类、推荐商品、玩法内容与经销合作入口；品牌页提供介绍、服务联系方式及安全信息入口。
\item {\bfseries 数据}：Banner 标题、图片和链接，推荐商品、品牌介绍和联系方式由后台读取。固定布局允许各模块在无内容时合理隐藏。
\item {\bfseries 规则}：推荐商品只包含已上架商品；安全说明引用可核实资料，不生成虚构认证或功效承诺。
\end{enumerate}

\begin{samepage}
\item {\bfseries 验收条件}：管理员改动 Banner 或推荐商品后，前台重新请求能看到新内容，无须重新编译代码。
\item {\bfseries 验证方法}：接口测试与业务场景验证；关联 TC-01。
\end{samepage}
\end{enumerate}

\reqheading{FR-03}{FR-03 商品列表、排序与分页}

\begin{enumerate}
\item {\bfseries 需求陈述}：系统应提供已上架商品的分页列表和稳定排序。
\item {\bfseries 业务理由}：帮助用户浏览商品集合并比较价格和适用性。
\item {\bfseries 适用条件}：所有访问者；列表只返回公开可售或可展示商品。
\item {\bfseries 行为与约束}：
\begin{enumerate}
\item {\bfseries 功能}：展示主图、名称、简述、适用年龄、零售价和有货／缺货状态；支持默认推荐顺序、价格升序和降序。
\item {\bfseries 规则}：前台只列出已上架商品；缺货商品可以浏览，明确提示不可购买。单页默认 12 项，列表接口最大单页数量为 50。
\item {\bfseries 交互}：展示总条数和当前页；更改排序或筛选时回到第一页；返回列表时保留当前查询条件。
\item {\bfseries 异常}：无商品显示空状态和返回／清除筛选入口；加载失败提供重试。
\end{enumerate}

\begin{samepage}
\item {\bfseries 验收条件}：有超过一页的数据时分页正确；页间无无故重复或遗漏；排序同值时采用固定次序，例如按商品 ID，保证结果稳定。
\item {\bfseries 验证方法}：接口测试与业务场景验证；关联 TC-02。
\end{samepage}
\end{enumerate}

\reqheading{FR-04}{FR-04 商品搜索与组合筛选}

\begin{enumerate}
\item {\bfseries 需求陈述}：系统应按关键词及组合条件查询商品，并保留查询状态。
\item {\bfseries 业务理由}：缩短用户寻找符合年龄、玩法和场景需求商品的过程。
\item {\bfseries 适用条件}：所有访问者；查询条件符合数据校验要求。
\item {\bfseries 行为与约束}：
\begin{enumerate}
\item {\bfseries 功能}：按商品名称或 SKU 关键词查询；按分类、适用年龄、玩法和室内／户外／两者皆可等使用场景筛选。
\item {\bfseries 规则}：关键词去除首尾空白，长度最多 100 个字符，英文匹配忽略大小写；空关键词视为不限制关键词。同一筛选维度本版本单选，不同维度取交集。
\item {\bfseries 场景语义}：选择室内时包含“室内”和“两者皆可”；选择户外时包含“户外”和“两者皆可”；选择“两者皆可”时只返回明确支持两种场景的商品。
\item {\bfseries 年龄语义}：用户输入期望适用年龄 \(a\)，匹配 \(age_{min}\leq a\leq age_{max}\) 的商品；未设置年龄上限时只比较下限。查询的是商品建议使用年龄，不保存儿童档案。
\item {\bfseries 交互与异常}：筛选状态保存在 URL 查询参数中，支持刷新与复制链接；无结果显示“未找到符合条件的商品”并可一键清除；非法参数提示或按接口约定拒绝。
\end{enumerate}

\begin{samepage}
\item {\bfseries 验收条件}：搜索结果满足关键词、年龄、场景和其他筛选条件；空关键词、无结果及非法页码按规则反馈。
\item {\bfseries 验证方法}：接口测试与业务场景验证；关联 TC-02、TC-03。
\end{samepage}
\end{enumerate}

\reqheading{FR-05}{FR-05 商品详情与购买入口}

\begin{enumerate}
\item {\bfseries 需求陈述}：系统应显示商品信息和与当前身份、可售状态一致的购买入口。
\item {\bfseries 业务理由}：使用户在购买或咨询前获得完整且一致的产品信息。
\item {\bfseries 适用条件}：所有访问者；加购需要启用的普通用户或经销商账户。
\item {\bfseries 行为与约束}：
\begin{enumerate}
\item {\bfseries 功能}：展示主图及附图、名称、SKU、零售价、库存状态、年龄范围、玩法、场景、材质、尺寸及单位、包装包含物、玩法说明、安全提示和公开资料下载。
\item {\bfseries 规则}：数量默认为 1，可输入正整数；可购买时提供加购入口，游客点击后进入登录并能返回原商品；已下架商品从列表消失，旧链接显示“已下架”且不能购买。
\item {\bfseries 价格边界}：公共详情始终显示零售价；有效经销商可通过明确标识的经销商入口查看专属价，避免将个人零售价格与经销商参考价格混用。
\item {\bfseries 异常}：图片缺失显示占位图；没有资料不显示空下载按钮；商品不存在返回未找到；缺货或下架后服务端也拒绝加购和付款。
\end{enumerate}

\begin{samepage}
\item {\bfseries 验收条件}：商品数据与后台一致，数量和购买限制生效，普通账户的响应数据不包含经销商价。
\item {\bfseries 验证方法}：接口测试与业务场景验证；关联 TC-04。
\end{samepage}
\end{enumerate}

\reqheading{FR-06}{FR-06 玩法文章与 FAQ}

\begin{enumerate}
\item {\bfseries 需求陈述}：系统应提供已发布玩法文章和 FAQ 的浏览能力。
\item {\bfseries 业务理由}：支持商品理解、自助使用和常见问题解答。
\item {\bfseries 适用条件}：所有访问者；内容须处于已发布状态。
\item {\bfseries 行为与约束}：
\begin{enumerate}
\item {\bfseries 功能}：浏览已发布文章列表和详情，查看分类 FAQ；文章可关联商品，FAQ 可关联具体产品；展示标题、正文和更新时间。
\item {\bfseries 规则}：草稿或下线内容不出现在公开列表，直接请求也不能读取；正文展示与编辑支持文字、图片、列表等基本内容。
\item {\bfseries 交互}：FAQ 可展开查看答案；文章提供返回列表入口；关联商品下架后不出现可购买的失效推荐。
\item {\bfseries 异常}：无内容时显示说明；正文中的危险脚本不得在浏览器执行。
\end{enumerate}

\begin{samepage}
\item {\bfseries 验收条件}：后台发布或下线内容后前台可见性随之改变，文章和 FAQ 读取后台维护的已发布内容。
\item {\bfseries 验证方法}：接口测试与业务场景验证；关联 TC-05。
\end{samepage}
\end{enumerate}

\reqheading{FR-07}{FR-07 下载中心与文件访问}

\begin{enumerate}
\item {\bfseries 需求陈述}：系统应按资料可见级别提供列表查询和文件下载。
\item {\bfseries 业务理由}：提供必要产品资料并保护经销商及内部文件。
\item {\bfseries 适用条件}：访问者满足文件对应权限；文件处于有效状态。
\item {\bfseries 行为与约束}：
\begin{enumerate}
\item {\bfseries 功能}：列出资料标题、类型、版本说明、适用商品和下载按钮；支持按资料类型或关联商品筛选。
\item {\bfseries 权限}：本版本文件可见级别为公开、经销商、内部。公开列表不返回非公开文件标题、路径或元数据；有效经销商可见公开和经销商资料；内部文件仅管理员可读。
\item {\bfseries 规则}：非公开文件每次下载由服务端校验权限，不能靠隐藏按钮或难猜的静态链接保护；资料可下线、替换当前版本，历史文件版本恢复不属于本版本范围。
\item {\bfseries 异常}：文件下线、资源不存在、权限失效或文件读取失败时提示对应错误，不暴露服务器存储路径。
\end{enumerate}

\begin{samepage}
\item {\bfseries 验收条件}：复制经销商下载地址到未登录会话中不能取到文件；在合作暂停后再次请求也不能下载。
\item {\bfseries 验证方法}：接口测试；关联 TC-06、TC-36。
\end{samepage}
\end{enumerate}

\reqheading{FR-08}{FR-08 经销商与购买渠道查询}

\begin{enumerate}
\item {\bfseries 需求陈述}：系统应按地区提供已发布的经销商和购买渠道信息。
\item {\bfseries 业务理由}：帮助用户找到公开销售或联系渠道。
\item {\bfseries 适用条件}：所有访问者；渠道经管理员明确发布。
\item {\bfseries 行为与约束}：
\begin{enumerate}
\item {\bfseries 功能}：按国家／地区和城市筛选公开经销渠道，展示名称、地址、联系方式及可选网站链接。
\item {\bfseries 规则}：只有管理员明确发布的渠道记录可见；企业审核通过不自动公开企业申请资料。公开渠道与申请资料分开管理。
\item {\bfseries 范围}：本版本以地区列表实现查找，不请求定位，不承诺“距离最近”，不依赖地图或地理编码服务。
\item {\bfseries 异常}：没有匹配渠道时，引导联系平台或使用站内模拟零售；没有网站时不显示空链接。
\end{enumerate}

\begin{samepage}
\item {\bfseries 验收条件}：未发布的渠道无法通过公开接口查询；筛选结果与后台维护的数据一致。
\item {\bfseries 验证方法}：接口测试与业务场景验证；关联 TC-07。
\end{samepage}
\end{enumerate}

\Needspace{12\baselineskip}
\subsubsection{账户与权限}\label{sec:322-账户与权限}

\reqheading{FR-09}{FR-09 成人账户注册}

\begin{enumerate}
\item {\bfseries 需求陈述}：系统应校验成人注册资料并建立普通用户账户。
\item {\bfseries 业务理由}：建立可识别的业务主体，支持个人记录与授权访问。
\item {\bfseries 适用条件}：未登录访问者提交注册资料和必要确认。
\item {\bfseries 行为与约束}：
\begin{enumerate}
\item {\bfseries 输入}：昵称、邮箱、密码、确认密码，已满 18 岁声明，以及使用说明和隐私说明确认。确认框默认不选中。
\item {\bfseries 处理}：校验通过后创建普通用户账户；邮箱在去除首尾空白并统一大小写后全局唯一。注册请求中的角色、用户状态等额外字段不能产生提权。
\item {\bfseries 规则}：昵称 2--30 字符；密码 8--64 字符，至少包含字母和数字；不对密码静默去空格或截断。注册不提供儿童出生日期等采集字段。
\item {\bfseries 异常}：邮箱格式不合法、重复邮箱、密码不一致、未确认成人身份或说明均拒绝注册，并提示需要修改的字段；失败不创建残缺账户。
\end{enumerate}

\begin{samepage}
\item {\bfseries 验收条件}：可用新账户登录；直接绕过前端提交非法注册数据仍被后端拒绝；管理员不能查看明文密码。
\item {\bfseries 验证方法}：接口测试与业务场景验证；关联 TC-08、TC-09。
\end{samepage}
\end{enumerate}

\reqheading{FR-10}{FR-10 登录、会话与退出}

\begin{enumerate}
\item {\bfseries 需求陈述}：系统应验证登录凭证，管理会话有效期，并允许用户退出。
\item {\bfseries 业务理由}：使账户访问可控制、可终止并具备基础滥用防护。
\item {\bfseries 适用条件}：登录适用于已注册账户；退出适用于有效会话。
\item {\bfseries 行为与约束}：
\begin{enumerate}
\item {\bfseries 功能}：邮箱密码登录；成功后进入原目标页或个人中心；退出后会话失效；会话过期时要求重新登录。
\item {\bfseries 规则}：同一登录入口可识别角色并引导至相应中心；错误邮箱或密码统一提示“邮箱或密码错误”；账户停用时拒绝访问受保护业务。
\item {\bfseries 项目阈值}：会话空闲 30 分钟失效；同一规范化邮箱在 10 分钟内连续失败 5 次后，暂停该邮箱密码尝试 10 分钟，反馈稍后重试。该限制不等同于管理员停用账户。
\item {\bfseries 异常}：登录失败不产生有效会话；退出后重放旧凭证应被拒绝；用户会话失效不清空已保存的购物车或历史订单。
\end{enumerate}

\begin{samepage}
\item {\bfseries 验收条件}：有效凭证建立会话；无效、受限或停用账户不能登录；超时与退出后的会话不能继续访问受保护资源。
\item {\bfseries 验证方法}：接口测试；关联 TC-10。
\end{samepage}
\end{enumerate}

\reqheading{FR-11}{FR-11 服务端权限与数据隔离}

\begin{enumerate}
\item {\bfseries 需求陈述}：系统应在每次受保护请求中校验身份、状态、角色和数据归属。
\item {\bfseries 业务理由}：防止未授权用户读取或修改他人订单、企业资料及平台数据。
\item {\bfseries 适用条件}：所有受保护接口及非公开文件访问。
\item {\bfseries 行为与约束}：
\begin{enumerate}
\item {\bfseries 功能}：按照第 2.3.3 节检查每一个受保护的读取和修改操作；用户只能查询或操作自己持有的业务记录。
\item {\bfseries 规则}：用户 ID、角色和经销商状态以服务端身份信息为准，不信任前端提交的 \code{userId}、\code{role} 或 \code{isDealer}；管理员入口隐藏与否不影响实际访问限制。
\item {\bfseries 失效处理}：账户停用后，已有会话的下一次受保护请求即失效；经销合作暂停后，下一次专属资源请求即被拒绝，不等待重新登录。
\item {\bfseries 异常}：未登录返回未认证；身份有效但权限不足返回无权限；访问他人私人记录统一按“记录不存在或不可访问”处理，避免泄露记录存在性。
\end{enumerate}

\begin{samepage}
\item {\bfseries 验收条件}：修改 URL 中订单 ID、直接请求后台接口、伪造角色字段以及使用暂停合作账户，均不能越权。
\item {\bfseries 验证方法}：接口测试与业务场景验证；关联 TC-11、TC-29。
\end{samepage}
\end{enumerate}

\reqheading{FR-12}{FR-12 个人资料与业务记录入口}

\begin{enumerate}
\item {\bfseries 需求陈述}：系统应提供个人资料维护及本人业务记录入口。
\item {\bfseries 业务理由}：使用户维护必要联系信息并集中跟踪业务进度。
\item {\bfseries 适用条件}：启用的普通用户或经销商账户，访问本人资料。
\item {\bfseries 行为与约束}：
\begin{enumerate}
\item {\bfseries 功能}：查看账户邮箱及身份状态，编辑昵称和可选联系电话；提供本人订单、咨询、经销申请和经销询价入口。
\item {\bfseries 范围}：邮箱作为登录标识，本版本不开放自助修改；地址在结算时填写，独立地址簿列入 P1；本版本不提供邮件找回密码。
\item {\bfseries 规则}：只有账户本人能修改个人资料；经销企业名称等审核资料不通过个人资料接口修改；刷新后保留已保存内容。
\item {\bfseries 异常}：昵称为空或过长、联系电话格式不符合第 3.4 节规则时拒绝保存，并保留用户输入以便修正。
\end{enumerate}

\begin{samepage}
\item {\bfseries 验收条件}：修改只影响本人，服务重启后可查到更新结果，普通用户不能通过资料接口赋予自己经销商或管理员身份。
\item {\bfseries 验证方法}：接口测试与业务场景验证；关联 TC-09、TC-12。
\end{samepage}
\end{enumerate}

\Needspace{12\baselineskip}
\subsubsection{零售交易}\label{sec:323-零售交易}

\reqheading{FR-13}{FR-13 登录购物车}

\begin{enumerate}
\item {\bfseries 需求陈述}：系统应为登录用户持久化保存购物车并校验商品和数量。
\item {\bfseries 业务理由}：支持用户汇总购买意向并在后续会话继续处理。
\item {\bfseries 适用条件}：启用的普通用户或经销商账户；加购商品符合可售条件。
\item {\bfseries 行为与约束}：
\begin{enumerate}
\item {\bfseries 功能}：登录用户添加商品、修改数量、删除商品和清空购物车；展示商品、单价、数量、行小计和总额。
\item {\bfseries 规则}：同一用户、同一商品只保留一条购物车项，重复加购累加数量；每行数量为 1--99 件，最多 20 个不同商品；购物车保存在服务端。
\item {\bfseries 价格与库存}：购物车读取当前零售价；加入和修改时检查上架状态与库存，但不预占库存。页面明确提示提交订单和付款时还会核验。
\item {\bfseries 异常}：商品后续缺货或下架时保留该行并标明无效，用户可移除；无效项未处理前不得继续整车结算。价格变动时提示用户查看更新金额。
\item {\bfseries 管理员限制}：管理员身份不能发起个人加购和结算；其订单处理能力按 FR-17、FR-18 执行。
\end{enumerate}

\begin{samepage}
\item {\bfseries 验收条件}：同账户重新登录仍能读取购物车，不同账户不共享；非法数量、超库存、重复商品和空车结算均正确处理。
\item {\bfseries 验证方法}：接口测试与业务场景验证；关联 TC-13、TC-14。
\end{samepage}
\end{enumerate}

\reqheading{FR-14}{FR-14 结算与待付款订单创建}

\begin{enumerate}
\item {\bfseries 需求陈述}：系统应依据服务端商品数据和用户确认内容创建唯一待模拟付款订单。
\item {\bfseries 业务理由}：形成可追踪的交易记录并固定成交内容。
\item {\bfseries 适用条件}：本人购物车非空且有效；收货信息和结算明细已确认。
\item {\bfseries 行为与约束}：
\begin{enumerate}
\item {\bfseries 输入}：收件人、联系电话、国家／地区、省市或相应地区信息、详细地址和可选订单备注；展示完整商品明细、运费 0 和模拟应付金额。
\item {\bfseries 处理}：本版本整车结算。服务端重新读取商品与价格、检查数量及库存，生成唯一订单号，保存订单、订单明细、地址和金额快照；状态为“待模拟付款”。
\item {\bfseries 价格变化}：结算预览至提交期间价格变化时，不静默按新价创建订单，应返回更新明细让用户再次确认；客户端传入的小计或合计不能作为计价依据。
\item {\bfseries 一致性}：创建订单及移除本次结算购物车项作为一次整体操作；失败不移除商品。每次确认使用业务请求标识，重试相同确认请求返回原订单，不重复创建。
\end{enumerate}

\begin{samepage}
\item {\bfseries 验收条件}：缺字段、不可售商品、库存不足、金额篡改、重复点击、数据库写入失败均不产生重复或不完整订单。待付款订单不预占库存。
\item {\bfseries 验证方法}：接口测试与业务场景验证；关联 TC-14、TC-15、TC-16、TC-17。
\end{samepage}
\end{enumerate}

\reqheading{FR-15}{FR-15 模拟付款}

\begin{enumerate}
\item {\bfseries 需求陈述}：系统应处理本人订单的模拟付款结果，并将成功付款与库存扣减一致保存。
\item {\bfseries 业务理由}：验证交易状态与库存处理，同时保证重复操作不产生额外影响。
\item {\bfseries 适用条件}：启用的普通用户或经销商账户，操作本人待模拟付款订单。
\item {\bfseries 行为与约束}：
\begin{enumerate}
\item {\bfseries 界面}：订单页明确写明“不发生真实扣款”，提供“模拟付款成功”和“模拟付款失败”两种验证动作，不填写银行卡或支付账户。
\item {\bfseries 成功处理}：只有本人待模拟付款订单可操作；沿用下单金额快照，不因后续改价重新计价。服务端再次验证账户、商品可售状态及库存，在同一事务中扣减所有订单项库存、记录模拟付款结果并转为待发货。
\item {\bfseries 失败处理}：模拟失败记录本次结果，订单仍为待模拟付款，不扣库存，用户可重试或取消；库存不足时逐项说明，并保留订单待用户处理。
\item {\bfseries 幂等性}：同一订单已经成功付款后再次提交成功或失败请求，不得再次扣库存、重复产生成功付款记录或退回待付款；返回当前状态或明确状态冲突。
\end{enumerate}

\begin{samepage}
\item {\bfseries 验收条件}：模拟失败和成功均产生对应的页面反馈和业务记录；两个订单争抢最后一件库存时至多一个成功；多商品订单任何一项不足时所有项均不扣减。
\item {\bfseries 验证方法}：接口测试与业务场景验证；关联 TC-17、TC-18、TC-19、TC-20。
\end{samepage}
\end{enumerate}

\reqheading{FR-16}{FR-16 本人订单列表与详情}

\begin{enumerate}
\item {\bfseries 需求陈述}：系统应向用户提供本人订单列表、快照明细和当前处理状态。
\item {\bfseries 业务理由}：使用户跟踪交易并保留与成交时一致的凭据。
\item {\bfseries 适用条件}：启用的普通用户或经销商账户，访问本人订单。
\item {\bfseries 行为与约束}：
\begin{enumerate}
\item {\bfseries 功能}：按创建时间倒序查看订单，可按状态筛选；详情展示订单号、创建时间、商品快照、地址快照、金额、模拟付款情况、状态历史和可执行动作。
\item {\bfseries 规则}：当前用户只能访问本人订单；商品改名、改价、下架后不改变历史订单；订单列表支持分页。
\item {\bfseries 入口}：待付款订单可继续模拟付款；可取消状态下显示取消；已发货订单可确认收货；订单详情可发起关联订单的售后咨询。
\item {\bfseries 异常}：无订单展示继续浏览入口；无权查看的订单不返回地址、姓名、金额等信息。
\end{enumerate}

\begin{samepage}
\item {\bfseries 验收条件}：直接修改订单 ID 不能读取他人信息；后台状态变化在刷新或重新请求后可见。
\item {\bfseries 验证方法}：接口测试；关联 TC-11、TC-21。
\end{samepage}
\end{enumerate}

\reqheading{FR-17}{FR-17 订单取消与库存恢复}

\begin{enumerate}
\item {\bfseries 需求陈述}：系统应在允许状态下取消订单，并按付款情况恢复库存和记录模拟退款。
\item {\bfseries 业务理由}：允许履约前撤销交易，同时保证库存与订单状态一致。
\item {\bfseries 适用条件}：订单属于本人或由管理员处理，且状态为待模拟付款或待发货。
\item {\bfseries 行为与约束}：
\begin{enumerate}
\item {\bfseries 功能}：用户可取消本人待模拟付款或待发货订单，需确认操作并填写或选择取消原因；管理员可在相同状态下代为取消并记录原因。
\item {\bfseries 待付款取消}：改为已取消，不改变库存。
\item {\bfseries 待发货取消}：整单取消，将本单已扣库存恢复一次，记录“模拟退款完成”；不发生真实退款，不保留可再次发货的状态。
\item {\bfseries 限制}：已发货、已完成不能直接取消，应提交售后咨询；已取消订单不可恢复或继续付款。付款、取消和发货须按当前状态串行校验并原子更新，付款后可以合法取消待发货订单，但取消与发货不能同时成功。
\end{enumerate}

\begin{samepage}
\item {\bfseries 验收条件}：重复取消不重复加库存；取消与发货同时执行不能同时成功；失效状态的操作返回可理解错误。
\item {\bfseries 验证方法}：接口测试；关联 TC-22、TC-23。
\end{samepage}
\end{enumerate}

\reqheading{FR-18}{FR-18 后台发货与用户确认收货}

\begin{enumerate}
\item {\bfseries 需求陈述}：系统应支持管理员记录整单模拟发货，并支持订单本人确认收货。
\item {\bfseries 业务理由}：形成可追踪的履约状态和完成结果。
\item {\bfseries 适用条件}：发货由管理员操作待发货订单；收货由本人操作已发货订单。
\item {\bfseries 行为与约束}：
\begin{enumerate}
\item {\bfseries 管理员操作}：在订单列表筛选待发货订单，核对订单和地址后填写模拟物流名称及运单号，确认整单发货，转为已发货并保存时间与操作人。
\item {\bfseries 用户操作}：本人已发货订单可确认收货，二次确认后转为已完成；重复确认不重复产生业务结果。
\item {\bfseries 规则}：本版本一次整单发货，不拆单、不部分发货、不调用真实物流；订单号、模拟交易标识和状态在前后台一致。
\item {\bfseries 限制}：管理员不能直接修改已创建订单的商品、数量、收货快照或金额；发生业务争议时按取消规则或售后工单流程处理。
\end{enumerate}

\begin{samepage}
\item {\bfseries 验收条件}：待付款、已取消或已完成订单不能发货；非本人不能确认收货；非法状态变化不改变数据。
\item {\bfseries 验证方法}：接口测试与业务场景验证；关联 TC-23、TC-24。
\end{samepage}
\end{enumerate}

\Needspace{12\baselineskip}
\subsubsection{咨询与售后支持}\label{sec:324-咨询与售后支持}

\reqheading{FR-19}{FR-19 联系、产品咨询与售后咨询提交}

\begin{enumerate}
\item {\bfseries 需求陈述}：系统应接收一般、产品和订单售后咨询并形成可查询工单。
\item {\bfseries 业务理由}：集中保存用户诉求并交由平台跟进。
\item {\bfseries 适用条件}：启用的普通用户或经销商账户；售后订单属于提交者。
\item {\bfseries 行为与约束}：
\begin{enumerate}
\item {\bfseries 输入}：咨询类型、主题、内容；产品咨询关联一个商品，售后咨询关联本人一个订单；联系方式默认使用账户邮箱，并允许填写可选电话。
\item {\bfseries 规则}：主题 2--100 字符，正文 10--2000 字符；订单关联必须由后端确认归属。本版本不开放用户图片或附件上传，不接收儿童照片或企业证照。
\item {\bfseries 处理}：提交后生成唯一工单编号，状态为待处理，保存提交者、类型、关联对象、内容和时间；页面显示编号及“查看我的咨询”入口。
\end{enumerate}

\begin{samepage}
\item {\bfseries 验收条件}：重复点击同一次提交不创建多条工单，后台可查到记录；无登录、非法关联、空白或超长内容被拒绝。售后咨询仅记录和跟进诉求，不等于已批准退货或退款。
\item {\bfseries 验证方法}：接口测试与业务场景验证；关联 TC-25、TC-26。
\end{samepage}
\end{enumerate}

\reqheading{FR-20}{FR-20 我的咨询与处理进度}

\begin{enumerate}
\item {\bfseries 需求陈述}：系统应允许用户查看、补充和关闭本人未关闭工单。
\item {\bfseries 业务理由}：使用户能够持续跟进服务处理并保留历史。
\item {\bfseries 适用条件}：启用的普通用户或经销商账户；工单归属当前用户。
\item {\bfseries 行为与约束}：
\begin{enumerate}
\item {\bfseries 功能}：查看本人咨询列表及详情，包含原始内容、关联商品／订单、处理状态、管理员公开回复及时间。
\item {\bfseries 规则}：用户可在工单未关闭时补充文本；工单关闭后只读，仍有新问题时提交新工单。补充内容不得覆盖历史对话。
\item {\bfseries 关闭}：用户可以关闭本人待处理、处理中或已回复工单，须确认；管理员也可关闭，需记录原因。
\item {\bfseries 异常}：访问其他用户工单、修改提交者或引用他人订单均拒绝；内部备注不能出现在用户响应中。
\end{enumerate}

\begin{samepage}
\item {\bfseries 验收条件}：用户能从提交追踪到回复或关闭结果；关闭后不能再追加消息或恢复状态。
\item {\bfseries 验证方法}：接口测试与业务场景验证；关联 TC-11、TC-25。
\end{samepage}
\end{enumerate}

\reqheading{FR-21}{FR-21 后台工单处理}

\begin{enumerate}
\item {\bfseries 需求陈述}：系统应支持管理员查询工单、公开回复、记录内部备注及关闭处理。
\item {\bfseries 业务理由}：使服务请求形成有结果、可追踪的处理过程。
\item {\bfseries 适用条件}：启用的管理员账户；写操作仅适用于未关闭工单。
\item {\bfseries 行为与约束}：
\begin{enumerate}
\item {\bfseries 功能}：管理员按类型、状态、时间查询工单，查看关联业务，标记处理中、添加公开回复、内部备注和关闭原因。
\item {\bfseries 状态}：新工单为待处理；开始处理后为处理中；发送公开回复后为已回复；已回复工单经用户补充后转为处理中，其余未关闭状态补充时保持原状态；关闭后只读。
\item {\bfseries 记录}：公开回复与内部备注分开保存和展示；记录操作人、时间和状态变化；不得修改用户原始提交内容。
\item {\bfseries 反馈方式}：公开回复保存后用户可在个人中心查看，本版本不要求邮件推送；空白回复不能将工单变为已回复。
\end{enumerate}

\begin{samepage}
\item {\bfseries 验收条件}：普通用户不能调用处理接口，处理结果可在用户端读取，关闭操作不删除工单及回复记录。
\item {\bfseries 验证方法}：接口测试与业务场景验证；关联 TC-25。
\end{samepage}
\end{enumerate}

\Needspace{12\baselineskip}
\subsubsection{经销合作}\label{sec:325-经销合作}

\reqheading{FR-22}{FR-22 经销商申请与本人申请记录}

\begin{enumerate}
\item {\bfseries 需求陈述}：系统应接收企业合作申请，限制重复申请并保存修订与审核历史。
\item {\bfseries 业务理由}：建立统一合作入口，提供稳定的审核依据。
\item {\bfseries 适用条件}：启用的普通用户账户，当前无待审核或已获批的合作资格。
\item {\bfseries 行为与约束}：
\begin{enumerate}
\item {\bfseries 输入}：企业名称、业务类型、国家／地区、城市、联系人、联系电话、合作邮箱、经营渠道说明、可选网站、合作意向及允许公开渠道信息的独立选择。
\item {\bfseries 规则}：只有已登录普通用户可申请；成人声明沿用账户规则。企业名称 2--100 字符，合作意向 10--2000 字符，业务类型从零售、批发、进口、教育／活动机构及其他中选择。
\item {\bfseries 提交}：生成申请编号并进入待审核；同一账户至多一份当前申请，待审核期间不可重复提交或修改，避免审核资料前后变化。
\item {\bfseries 重提}：被驳回后可修订重新提交，保留上一次审核原因与提交版本；已通过或暂停合作的用户不能重新申请以绕过合作状态。
\end{enumerate}

\begin{samepage}
\item {\bfseries 验收条件}：用户能查看自己的申请内容和状态；企业名与地区相同的申请向管理员标记疑似重复，由管理员判断，不仅凭重名自动合并企业。
\item {\bfseries 验证方法}：接口测试与业务场景验证；关联 TC-27、TC-28。
\end{samepage}
\end{enumerate}

\reqheading{FR-23}{FR-23 申请审核与经销权限生效}

\begin{enumerate}
\item {\bfseries 需求陈述}：系统应允许管理员审核当前申请版本，并一致更新审核结果、企业记录和经销资格。
\item {\bfseries 业务理由}：确保合作审批结果与实际访问权限一致。
\item {\bfseries 适用条件}：启用的管理员账户；申请处于待审核且版本有效。
\item {\bfseries 行为与约束}：
\begin{enumerate}
\item {\bfseries 功能}：管理员查看待审核申请及历史，选择通过或驳回；驳回原因必填，并与内部备注分开。
\item {\bfseries 通过处理}：在一次事务中更新申请状态、建立或关联唯一企业记录并激活该用户的经销资格；仅可关联已归属于当前申请用户的企业记录，不得占用其他账户的主账户关系。重复审核不能创建重复企业或多份资格。
\item {\bfseries 结果反馈}：用户可在个人中心查看审核结果；通过后刷新即可进入经销商中心；被驳回只具备普通用户能力。
\item {\bfseries 边界}：申请通过不自动创建公开门店，不自动发布申请人的邮箱或电话；同名企业由管理员判断，不自动合并企业、共享资格或覆盖原有主账户。两个管理员并发给出不同结论时仅第一次有效处理成功。
\end{enumerate}

\begin{samepage}
\item {\bfseries 验收条件}：非管理员不能审核；审核后权限与显示结果一致，不存在“已通过但无企业记录”或“被驳回但可访问专属价”的部分成功。
\item {\bfseries 验证方法}：接口测试与业务场景验证；关联 TC-27、TC-28。
\end{samepage}
\end{enumerate}

\reqheading{FR-24}{FR-24 经销商目录、价格与资料}

\begin{enumerate}
\item {\bfseries 需求陈述}：系统应仅向有效经销商提供专属目录、参考价和经销资料。
\item {\bfseries 业务理由}：支持采购准备并隔离未授权的商业信息。
\item {\bfseries 适用条件}：启用账户且经销合作有效；商品启用经销业务。
\item {\bfseries 行为与约束}：
\begin{enumerate}
\item {\bfseries 功能}：有效经销商可浏览已上架且启用经销业务的商品，查看 SKU、统一经销商单价、最小询价数量、现有库存及管理员维护的参考交期，下载经销商资料。
\item {\bfseries 规则}：经销商价必须单独标注币种与“参考单价，实际合作以回复为准”；参考交期为人工维护信息，不表示外部仓储系统已经确认供货。未配置经销商价的商品不进入专属目录。
\item {\bfseries 权限}：公共商品接口不携带专属价；经销接口检查合作状态；本版本所有有效经销商共享目录和价格，不实现企业专属定价或不同授权商品集。
\item {\bfseries 与零售的关系}：经销商仍可在公开商品页按零售价进行个人模拟购买；经销商参考价不能通过参数提交到零售订单中使用。
\end{enumerate}

\begin{samepage}
\item {\bfseries 验收条件}：同一商品在两个入口显示清楚的价格类型，普通用户和暂停合作用户不能获取专属响应或下载文件。
\item {\bfseries 验证方法}：接口测试与业务场景验证；关联 TC-29、TC-32。
\end{samepage}
\end{enumerate}

\reqheading{FR-25}{FR-25 经销商询价与回复}

\begin{enumerate}
\item {\bfseries 需求陈述}：系统应处理经销商的多商品询价、管理员回复和询价关闭。
\item {\bfseries 业务理由}：使采购需求与供货沟通形成可追踪记录。
\item {\bfseries 适用条件}：有效经销商可新建；本人可读取和关闭历史询价；管理员可处理。
\item {\bfseries 行为与约束}：
\begin{enumerate}
\item {\bfseries 输入}：必须从专属目录选择 1--20 个不同商品并填写各行数量；期望交付日期、交付说明、用途和备注均为选填。交付说明、用途和备注各不超过 2000 字符。数量为正整数，不低于该商品最小询价数量，上限 9999。
\item {\bfseries 处理}：提交时校验资格、商品、重复行和数量，保存企业、提交者、商品名称／SKU／参考价快照和需求，生成询价编号，状态为待处理。
\item {\bfseries 后台回复}：管理员可标记处理中，填写回复文字，可附每行参考单价、预计交期说明；保存后为已回复。经销商可查看回复，并可关闭本人任一未关闭询价，管理员也可说明原因后关闭；关闭后只读。
\item {\bfseries 范围}：询价数量允许超过当前库存，以便沟通供货；提交和回复均不扣库存、不预占、不形成付款义务。需要改动已提交的采购明细时关闭原询价并新建，保留历史。
\item {\bfseries 权限与异常}：只允许本人企业查看询价；暂停合作后可读已有询价和关闭本人询价，不能新建或获取新的专属目录信息。原记录中已获得的价格快照保留可读。
\end{enumerate}

\begin{samepage}
\item {\bfseries 验收条件}：后台能回复、用户能读到结果；数量不足、非目录商品、跨企业查询及重复请求均有明确处理；回复不得触发订单创建、付款或库存扣减。
\item {\bfseries 验证方法}：接口测试与业务场景验证；关联 TC-30、TC-31、TC-32。
\end{samepage}
\end{enumerate}

\Needspace{12\baselineskip}
\subsubsection{运营管理}\label{sec:326-运营管理}

\reqheading{FR-26}{FR-26 商品建立、编辑与上下架}

\begin{enumerate}
\item {\bfseries 需求陈述}：系统应支持商品草稿建立、编辑、发布、下架和重新上架。
\item {\bfseries 业务理由}：使商品信息由运营维护，并控制对外展示与交易时点。
\item {\bfseries 适用条件}：启用的管理员账户；发布及已上架编辑满足字段要求。
\item {\bfseries 行为与约束}：
\begin{enumerate}
\item {\bfseries 功能}：管理员新建商品草稿，维护第 3.4 节要求的商品资料、图片、零售价、经销商价及关联文件，可发布、编辑和下架。初始库存在建档时录入，未录入则为 0；建档后的库存变更只能通过 FR-27 执行，普通商品编辑不得直接覆盖库存。
\item {\bfseries 发布条件}：名称、唯一 SKU、分类、简述、年龄下限、玩法、场景、规格与单位、包装内容、安全提示、主图、有效零售价和非负整数库存齐全。启用经销业务时还需经销商价、最小询价数量及参考交期说明。
\item {\bfseries 规则}：草稿可以保存不完整字段；发布时必须完整校验；已上架商品编辑后仍须满足发布条件。SKU 可以在草稿阶段补录，但首次赋值后不得修改，且须满足唯一性要求。
\item {\bfseries 保留数据}：有订单或询价引用的商品不物理删除，本版本统一以下架停止展示和交易；历史订单、询价保留当时快照。
\end{enumerate}

\begin{samepage}
\item {\bfseries 验收条件}：完成草稿、发布、编辑、下架、重新上架流程；非法价格、重复 SKU、遗漏主图或分类不能发布；下架立即影响后续加购和付款校验。
\item {\bfseries 验证方法}：接口测试与业务场景验证；关联 TC-04、TC-33、TC-34。
\end{samepage}
\end{enumerate}

\reqheading{FR-27}{FR-27 分类、价格与库存维护}

\begin{enumerate}
\item {\bfseries 需求陈述}：系统应支持分类、两类价格和库存维护，并保护关联记录与并发更新。
\item {\bfseries 业务理由}：保持商品组织、价格来源和可售库存一致。
\item {\bfseries 适用条件}：启用的管理员账户；变更满足唯一性、引用和余额约束。
\item {\bfseries 行为与约束}：
\begin{enumerate}
\item {\bfseries 分类}：管理员维护单层分类名称、说明、顺序和启用状态；名称去除首尾空白后唯一。被商品引用的分类不能删除或停用，需先迁移商品；前台只显示启用分类。
\item {\bfseries 玩法与场景}：使用固定枚举或受控标签，不能因自由输入“室内”“室 内”等重复词造成筛选失效。
\item {\bfseries 价格}：零售与经销价格分别保存，金额采用精确小数或最小货币单位整数；修改价格不改变已创建订单与询价快照。
\item {\bfseries 库存}：后台采用“增加／减少数量＋调整原因”的操作，原子更新库存，不用旧页面上的绝对库存值覆盖并发扣减；减少后不得为负；系统付款扣减和取消返还也形成库存变动记录。
\end{enumerate}

\begin{samepage}
\item {\bfseries 验收条件}：前台下一次读取能看到新价格和库存；后台调整与用户付款同时发生时记录可核对、不丢失扣减；未获授权者不能调整。
\item {\bfseries 验证方法}：接口测试；关联 TC-33、TC-34。
\end{samepage}
\end{enumerate}

\reqheading{FR-28}{FR-28 文章、FAQ、Banner 与品牌信息维护}

\begin{enumerate}
\item {\bfseries 需求陈述}：系统应支持文章、FAQ、Banner 和品牌信息的编辑、发布及下线。
\item {\bfseries 业务理由}：使公开内容与品牌运营信息无需代码发布即可更新。
\item {\bfseries 适用条件}：启用的管理员账户；发布内容满足必填和安全要求。
\item {\bfseries 行为与约束}：
\begin{enumerate}
\item {\bfseries 功能}：维护文章、FAQ 和品牌介绍；配置首页 Banner 图片、标题、按钮文字与目标页面，选择推荐商品并设置顺序；维护联系方式。
\item {\bfseries 状态}：文章、FAQ、Banner 支持草稿／未启用、发布／启用、下线；发布前校验标题和必要内容；Banner 必须包含有效图片及目标链接。
\item {\bfseries 基础 SEO}：公开商品和文章提供独立 URL 和明确页面标题；可填写页面描述，未填时使用简述。复杂站点地图和结构化数据不属于本版本强制范围。
\item {\bfseries 范围}：不要求拖拽任意模块、定时发布、多语言版本或版本回滚；本版本采用固定模板与编辑表单。
\end{enumerate}

\begin{samepage}
\item {\bfseries 验收条件}：管理员通过编辑界面完成内容维护；未发布数据不外泄，危险 HTML 不执行，后台修改无需重新部署前端。
\item {\bfseries 验证方法}：接口测试与业务场景验证；关联 TC-05、TC-35。
\end{samepage}
\end{enumerate}

\reqheading{FR-29}{FR-29 媒体和下载文件管理}

\begin{enumerate}
\item {\bfseries 需求陈述}：系统应安全接收管理员素材上传，管理其属性、权限、引用和替换。
\item {\bfseries 业务理由}：保证公开内容可用，并防止私有文件泄漏或无效引用。
\item {\bfseries 适用条件}：启用的管理员账户；文件符合类型、大小和引用约束。
\item {\bfseries 行为与约束}：
\begin{enumerate}
\item {\bfseries 功能}：管理员上传商品／内容图片及 PDF 下载文件，填写标题、类型、关联商品、版本说明和可见级别，支持替换和下线。
\item {\bfseries 限制}：图片仅支持 JPEG、PNG、WebP，单个不超过 5 MiB；下载资料仅支持 PDF，单个不超过 10 MiB。后端同时校验扩展名、声明类型及实际可识别文件内容。
\item {\bfseries 规则}：文件名经安全处理，存储标识由系统生成；文件不可作为服务器代码执行。经销商和内部文件不放入可直接公开访问的静态目录。
\item {\bfseries 引用关系}：已被上架商品或发布内容使用的图片，移除前应更换引用或先下线引用方；更换资料后旧下载标识失效，旧文件不再由业务入口访问。
\end{enumerate}

\begin{samepage}
\item {\bfseries 验收条件}：超限、类型伪装、未经授权上传均拒绝；上传失败不保留可见的空文件记录；相同名称文件不覆盖其他资源。
\item {\bfseries 验证方法}：接口测试与业务场景验证；关联 TC-06、TC-35、TC-36。
\end{samepage}
\end{enumerate}

\reqheading{FR-30}{FR-30 经销商企业与公开渠道维护}

\begin{enumerate}
\item {\bfseries 需求陈述}：系统应支持企业合作状态维护和独立的公开渠道发布。
\item {\bfseries 业务理由}：使合作资格、公开展示和内部资料保持清楚边界。
\item {\bfseries 适用条件}：启用的管理员账户；关键状态变更需说明原因。
\item {\bfseries 行为与约束}：
\begin{enumerate}
\item {\bfseries 企业管理}：管理员查询企业资料和来源申请，可添加内部备注、暂停或恢复合作，操作需原因；相关历史记录保留。
\item {\bfseries 资料更新}：企业关键资料由管理员根据咨询记录修正并留痕；本版本不开放企业用户自行修改审核资料。
\item {\bfseries 公开渠道管理}：独立维护可公开的名称、地区、地址、电话和网站，设置发布或下线；使用申请数据创建渠道时检查公开意愿并选择公开字段。
\item {\bfseries 联动}：与企业关联的公开渠道在合作暂停时自动下线，恢复合作后由管理员重新发布；独立独立渠道由管理员控制状态。
\end{enumerate}

\begin{samepage}
\item {\bfseries 验收条件}：暂停立即撤销专属能力，恢复可恢复能力；未经发布的申请资料不会自动进入公开查询。
\item {\bfseries 验证方法}：接口测试与业务场景验证；关联 TC-07、TC-32。
\end{samepage}
\end{enumerate}

\reqheading{FR-31}{FR-31 用户账户状态管理}

\begin{enumerate}
\item {\bfseries 需求陈述}：系统应支持管理员停用和恢复用户账户，并及时使访问控制生效。
\item {\bfseries 业务理由}：控制异常账户访问且保留既有业务记录。
\item {\bfseries 适用条件}：启用的管理员账户；目标为普通用户或经销商账户。
\item {\bfseries 行为与约束}：
\begin{enumerate}
\item {\bfseries 功能}：管理员按邮箱、昵称、角色或状态查询用户，查看必要账户信息和关联业务，停用或恢复普通用户／经销商账户。
\item {\bfseries 规则}：停用和恢复需记录原因及操作人；停用影响全部受保护业务且使旧会话不可继续使用，保留订单、工单、企业和审核记录；管理员仍可继续处理已有业务。
\item {\bfseries 保护}：本版本不允许通过该模块停用管理员自己或修改管理员角色；不得显示密码、密码摘要、会话凭证或其他秘密值。
\item {\bfseries 范围}：不做用户数据批量导出、在线创建管理员或不可恢复的账户删除；恢复账户不自动恢复此前被暂停的经销合作状态。
\end{enumerate}

\begin{samepage}
\item {\bfseries 验收条件}：普通用户不能调用账户管理接口；旧会话在账户停用后即被拒绝；恢复后原业务记录仍可查询。
\item {\bfseries 验证方法}：接口测试；关联 TC-10、TC-32。
\end{samepage}
\end{enumerate}

\reqheading{FR-32}{FR-32 后台总览与关键操作记录}

\begin{enumerate}
\item {\bfseries 需求陈述}：系统应展示定义明确的业务待办与统计，并保存关键操作记录。
\item {\bfseries 业务理由}：支持运营处理、数据核对和问题追踪。
\item {\bfseries 适用条件}：启用的管理员账户；统计基于已保存业务数据。
\item {\bfseries 行为与约束}：
\begin{enumerate}
\item {\bfseries 总览}：显示已上架商品数、启用用户数、待审核申请数、待处理工单／询价数、待发货订单数，以及选定时间段内创建订单数和模拟净成交金额。
\item {\bfseries 统计口径}：工单和询价的待办包含待处理与处理中；订单创建数含所有状态；模拟净成交金额按模拟付款成功时间归入时间段，累计成功金额减去这些订单当前已完成的模拟退款金额。取消后重查历史时间段，金额会相应减少。
\item {\bfseries 时间规则}：数据库时间统一保存，展示按 Asia/Shanghai；时间段采用起点包含、终点不包含。无数据时显示 0，不用虚构图表数据。
\item {\bfseries 操作记录}：记录商品上下架与改价、库存调整、订单状态变化、申请审核、合作暂停／恢复、账户停用／恢复、文件权限变化；至少含操作人、时间、对象、动作、结果和必要原因。
\item {\bfseries 保密与保留}：日志不含明文密码、会话凭证或完整非必要个人数据；后台日志只读，不提供普通删除功能。
\end{enumerate}

\begin{samepage}
\item {\bfseries 验收条件}：待办卡片可进入相应列表；统计可与测试订单核对；关键操作与对应业务变更一致，失败操作不得记为业务成功。
\item {\bfseries 验证方法}：接口测试与业务场景验证；关联 TC-37。
\end{samepage}
\end{enumerate}

\Needspace{10\baselineskip}
\subsection{业务规则与状态约束}\label{sec:33-业务规则与状态约束}

本节规则约束所有相关功能。状态表只列出允许变化；未列出的状态变更应拒绝，重复请求处理按 BR-06 执行。

\reqheading[3]{BR-01}{BR-01 商品价格和金额}

零售订单金额由服务端计算。设订单有 \(n\) 种商品，第 \(i\) 种商品下单时零售单价为 \(p_i\)，数量为 \(q_i\)：

\[
M=\sum_{i=1}^{n}p_iq_i+S+T-D
\]

本版本运费 \(S=0\)，额外税费 \(T=0\)，优惠 \(D=0\)。金额精确到分，存储和计算不得引入浮点舍入导致的分币误差。

\begin{enumerate}
\item 购物车使用当前价；结算创建订单时生成成交快照；订单创建后的改价不改变该待付款订单及历史订单金额。
\item 本版本不设置待付款自动过期；用户可继续付款或取消，付款时仍检查商品是否可售、库存是否足够。管理员可按 FR-17 取消长期未处理的订单。
\item 经销商参考价不参与零售计价；询价回复不构成价格锁定订单。
\item 已创建订单不允许直接改价。需要改动时，在允许状态取消后重新下单。
\end{enumerate}

\reqheading[3]{BR-02}{BR-02 库存变化}

\begingroup
\smalltable
\begin{longtable}{>{\raggedright\arraybackslash}p{\dimexpr 0.290\linewidth-2\tabcolsep\relax}>{\raggedright\arraybackslash}p{\dimexpr 0.240\linewidth-2\tabcolsep\relax}>{\raggedright\arraybackslash}p{\dimexpr 0.470\linewidth-2\tabcolsep\relax}}
\caption{库存变化规则}\label{tab:12}\\
\toprule
{\bfseries 动作} & {\bfseries 库存变化} & {\bfseries 必要条件} \\
\midrule
\endfirsthead
\multicolumn{3}{r}{\footnotesize 表~\thetable~（续）}\\
\toprule
{\bfseries 动作} & {\bfseries 库存变化} & {\bfseries 必要条件} \\
\midrule
\endhead
\midrule
\multicolumn{3}{r}{\footnotesize 续下页}\\
\endfoot
\bottomrule
\endlastfoot
加入购物车／提交待付款订单 & 不变 & 当时商品可售、数量合法且不超过当前库存。 \\
\tabledivider
模拟付款失败 & 不变 & 订单仍为待模拟付款。 \\
\tabledivider
模拟付款成功 & 减少订单数量 & 在一次事务中检查并扣减全部商品，转为待发货。 \\
\tabledivider
取消待付款订单 & 不变 & 尚未成功付款。 \\
\tabledivider
取消待发货订单 & 增加此前扣减数量 & 成功取消一次，仅返还一次。 \\
\tabledivider
发货／确认收货 & 不变 & 已在付款成功时扣减。 \\
\tabledivider
经销询价／工单提交和回复 & 不变 & 没有交易扣减含义。 \\
\tabledivider
后台调整 & 按输入增减 & 调整原因必填，结果不小于 0。 \\
\end{longtable}
\endgroup

对任意商品，始终要求库存 \(stock\geq0\)。库存校验与更新不可分成两个互不保护的步骤；并发购买同一库存时，成功订单的累计扣减量不得超过可用库存；数据库并发控制方式由设计阶段确定。

\reqheading[3]{BR-03}{BR-03 零售订单状态机}

\begingroup
\smalltable
\begin{longtable}{>{\raggedright\arraybackslash}p{\dimexpr 0.260\linewidth-2\tabcolsep\relax}>{\raggedright\arraybackslash}p{\dimexpr 0.260\linewidth-2\tabcolsep\relax}>{\raggedright\arraybackslash}p{\dimexpr 0.150\linewidth-2\tabcolsep\relax}>{\raggedright\arraybackslash}p{\dimexpr 0.330\linewidth-2\tabcolsep\relax}}
\caption{零售订单状态机}\label{tab:13}\\
\toprule
{\bfseries 当前状态} & {\bfseries 允许动作} & {\bfseries 执行者} & {\bfseries 下一状态与结果} \\
\midrule
\endfirsthead
\multicolumn{4}{r}{\footnotesize 表~\thetable~（续）}\\
\toprule
{\bfseries 当前状态} & {\bfseries 允许动作} & {\bfseries 执行者} & {\bfseries 下一状态与结果} \\
\midrule
\endhead
\midrule
\multicolumn{4}{r}{\footnotesize 续下页}\\
\endfoot
\bottomrule
\endlastfoot
待模拟付款 \code{PENDING\_\allowbreak{}PAYMENT} & 模拟成功且库存校验通过 & 本人 & 待发货 \code{PAID}，扣库存。 \\
\tabledivider
待模拟付款 & 模拟失败／库存不足 & 本人 & 仍待模拟付款，不扣库存。 \\
\tabledivider
待模拟付款 & 取消 & 本人或\newline 管理员 & 已取消 \code{CANCELLED}，库存不变。 \\
\tabledivider
待发货 \code{PAID} & 模拟发货 & 管理员 & 已发货 \code{SHIPPED}。 \\
\tabledivider
待发货 & 取消并完成模拟退款 & 本人或\newline 管理员 & 已取消，库存恢复一次。 \\
\tabledivider
已发货 \code{SHIPPED} & 确认收货 & 本人 & 已完成 \code{COMPLETED}。 \\
\tabledivider
已完成／已取消 & 查看／提交咨询 & 本人，\newline 管理员\newline 可查看 & 订单状态不变。 \\
\end{longtable}
\endgroup

付款记录可区分失败、成功、已模拟退款；其结果与订单状态保持一致。退款说明只适用于已模拟付款且未发货的整单取消。已发货后的真实退换货诉求在 P0 中仅形成售后工单，执行退款能力属于 P1。

\reqheading[3]{BR-04}{BR-04 经销商申请与合作状态}

\begingroup
\smalltable
\begin{longtable}{>{\raggedright\arraybackslash}p{\dimexpr 0.130\linewidth-2\tabcolsep\relax}>{\raggedright\arraybackslash}p{\dimexpr 0.430\linewidth-2\tabcolsep\relax}>{\raggedright\arraybackslash}p{\dimexpr 0.440\linewidth-2\tabcolsep\relax}}
\caption{申请与合作状态}\label{tab:14}\\
\toprule
{\bfseries 对象} & {\bfseries 状态与允许变化} & {\bfseries 约束} \\
\midrule
\endfirsthead
\multicolumn{3}{r}{\footnotesize 表~\thetable~（续）}\\
\toprule
{\bfseries 对象} & {\bfseries 状态与允许变化} & {\bfseries 约束} \\
\midrule
\endhead
\midrule
\multicolumn{3}{r}{\footnotesize 续下页}\\
\endfoot
\bottomrule
\endlastfoot
申请 & 未提交 $\rightarrow$ 待审核 \code{PENDING} & 未提交是表单阶段，本版本不要求服务端草稿。 \\
\tabledivider
申请 & 待审核 $\rightarrow$ 通过 \code{APPROVED}／驳回 \code{REJECTED} & 仅管理员；处理过的申请不能再次直接审核。 \\
\tabledivider
申请 & 驳回 $\rightarrow$ 修订重新待审核 & 本人操作，保留旧提交版本与结果。 \\
\tabledivider
合作 & 有效 \code{ACTIVE} $\rightarrow$ 暂停 \code{SUSPENDED} & 仅管理员；不回写旧申请为驳回。 \\
\tabledivider
合作 & 暂停 $\rightarrow$ 有效 & 仅管理员；不重复创建企业。 \\
\tabledivider
账户 & 启用／停用 & 独立于申请与合作状态；账户停用优先阻止全部受保护访问。 \\
\end{longtable}
\endgroup

\reqheading[3]{BR-05}{BR-05 工单和询价状态}

\begingroup
\smalltable
\begin{longtable}{>{\raggedright\arraybackslash}p{\dimexpr 0.160\linewidth-2\tabcolsep\relax}>{\raggedright\arraybackslash}p{\dimexpr 0.400\linewidth-2\tabcolsep\relax}>{\raggedright\arraybackslash}p{\dimexpr 0.440\linewidth-2\tabcolsep\relax}}
\caption{工单和询价状态}\label{tab:15}\\
\toprule
{\bfseries 对象} & {\bfseries 变化} & {\bfseries 说明} \\
\midrule
\endfirsthead
\multicolumn{3}{r}{\footnotesize 表~\thetable~（续）}\\
\toprule
{\bfseries 对象} & {\bfseries 变化} & {\bfseries 说明} \\
\midrule
\endhead
\midrule
\multicolumn{3}{r}{\footnotesize 续下页}\\
\endfoot
\bottomrule
\endlastfoot
工单和询价 & 提交 $\rightarrow$ 待处理 \code{NEW} & 记录唯一编号、发起人和关联业务。 \\
\tabledivider
工单和询价 & 待处理 $\rightarrow$ 处理中 \code{PROCESSING} & 管理员开始处理。 \\
\tabledivider
工单和询价 & 待处理／处理中 $\rightarrow$ 已回复 \code{REPLIED} & 保存非空公开回复；管理员可继续追加回复，状态仍为已回复。 \\
\tabledivider
工单 & 已回复 $\rightarrow$ 处理中 & 用户补充问题；待处理阶段补充仍保持待处理，处理中补充仍保持处理中。 \\
\tabledivider
工单和询价 & 待处理／处理中／已回复 $\rightarrow$ 已关闭 \code{CLOSED} & 本人或管理员关闭，保留原始内容与历史。 \\
\end{longtable}
\endgroup

询价不在原单追加采购明细或发起来回改价；本版本用户只读回复并关闭，需要新需求时另建询价。暂停经销合作但账户未停用的用户仍能查看、关闭历史询价。

\reqheading[3]{BR-06}{BR-06 重复请求、并发与失败恢复}

\begin{enumerate}
\item 创建订单、工单、申请和询价时，同一次提交的业务请求标识应绑定当前用户、动作和请求内容；相同标识及内容重试返回既有结果，相同标识但内容不同应拒绝。
\item 订单创建请求的去重结果至少保留 24 小时；订单付款和取消等状态动作在订单保留期间持续受状态限制，不因去重缓存到期再次发生资金或库存效果。
\item 申请版本与重复审核检查同时执行；审核时校验当前申请版本，避免旧页面审批新提交的资料。
\item 付款、取消、发货、库存调整采用一致的并发控制，确保任何成功结果都符合状态机；并发冲突返回当前状态供重新确认。
\item 网络超时可能代表请求尚未返回，也可能已经成功。页面应提示“结果暂未确认，可查询记录或重试”，而不是直接断言业务失败。
\item 业务变更失败时不得留下半成品，例如只扣一部分库存、删除购物车却未生成订单、审核通过却未授予资格。业务状态和必需的状态历史在同一成功边界保存。
\end{enumerate}

\reqheading[3]{BR-07}{BR-07 数据保留与可见性}

商品、订单、企业、申请、询价、工单及关键操作记录保留关联关系。前台下架、合作暂停、账户停用和关闭工单不等同于删除历史数据。快照负责保存当时事实，当前主数据负责决定下一次操作是否允许，两者不能混用。

公开渠道信息、用户个人联系信息和经销申请信息应使用不同的读取边界；公开接口采用允许返回字段清单，不能先返回所有字段再由前端隐藏。

\Needspace{10\baselineskip}
\subsection{数据要求}\label{sec:34-数据要求}

\Needspace{12\baselineskip}
\subsubsection{核心业务对象}\label{sec:341-核心业务对象}

本节给出需要保存的业务信息，不要求开发阶段机械对应为同名数据库表。

\begingroup
\smalltable
\begin{longtable}{>{\raggedright\arraybackslash}p{\dimexpr 0.170\linewidth-2\tabcolsep\relax}>{\raggedright\arraybackslash}p{\dimexpr 0.470\linewidth-2\tabcolsep\relax}>{\raggedright\arraybackslash}p{\dimexpr 0.360\linewidth-2\tabcolsep\relax}}
\caption{核心业务对象}\label{tab:16}\\
\toprule
{\bfseries 对象} & {\bfseries 关键字段} & {\bfseries 关联与约束} \\
\midrule
\endfirsthead
\multicolumn{3}{r}{\footnotesize 表~\thetable~（续）}\\
\toprule
{\bfseries 对象} & {\bfseries 关键字段} & {\bfseries 关联与约束} \\
\midrule
\endhead
\midrule
\multicolumn{3}{r}{\footnotesize 续下页}\\
\endfoot
\bottomrule
\endlastfoot
用户 & ID、规范化邮箱、密码摘要、昵称、电话、状态、成人确认时间、说明版本与确认时间 & 邮箱唯一；密码摘要不对接口返回。 \\
\tabledivider
会话 & 用户、会话标识、失效时间／最后活动时间、撤销状态 & 退出、停用和超时均可失效。 \\
\tabledivider
分类 & ID、名称、排序、状态 & 商品引用中的分类不能删除或停用。 \\
\tabledivider
商品 & ID、SKU、名称、分类、描述、年龄区间、玩法、场景、规格、图片、状态、零售价、经销业务启用标记、经销价、库存、最小询价数量、参考交期 & SKU 唯一；一商品一 SKU；公开响应剔除经销商字段。 \\
\tabledivider
购物车项 & 用户、商品、数量、更新时间 & 用户和商品联合唯一。 \\
\tabledivider
订单及明细 & 订单号、用户、状态、时间、商品／SKU／数量／单价快照、地址快照、金额、模拟发货信息 & 订单号唯一；归属不可由普通用户修改。 \\
\tabledivider
付款与库存记录 & 订单、动作、数量／金额、结果、时间、原因、业务请求标识 & 成功付款及返还只产生一次有效业务效果。 \\
\tabledivider
工单及消息 & 编号、用户、类型、关联商品／订单、正文、状态、公开回复／内部备注、时间 & 售后关联订单必须属于提交者。 \\
\tabledivider
经销申请及版本 & 编号、用户、企业与联系人信息、版本、状态、审核人、公开原因、内部备注、公开意愿 & 一个用户至多一份当前申请；旧版本不可覆盖。 \\
\tabledivider
经销企业 & ID、主用户、来源申请、企业资料、合作状态、内部备注 & 主用户唯一，不做子账户。 \\
\tabledivider
公开渠道 & 名称、地区、地址、公开联系方式、发布状态、可选企业关联 & 与申请资料分开；关联企业暂停时下线。 \\
\tabledivider
询价及明细／回复 & 编号、企业、发起用户、商品快照、数量、需求、状态、回复、参考报价与交期、时间 & 仅本人企业及管理员访问，不视为正式订单。 \\
\tabledivider
文章／FAQ／Banner／品牌设置 & 标题、正文或图片、链接、关联商品、排序、状态、页面描述 & 仅发布内容进入公开查询。 \\
\tabledivider
文件 & ID、显示名称、安全存储标识、实际类型、大小、权限、版本说明、关联对象、状态 & 不将服务器路径或私有标识暴露给未授权角色。 \\
\tabledivider
关键操作记录 & 操作人、动作、对象、时间、结果、原因及必要变化摘要 & 只读保留；不存储秘密凭证。 \\
\end{longtable}
\endgroup

\Needspace{12\baselineskip}
\subsubsection{数据字典与输入约束}\label{sec:342-数据字典与输入约束}

所有限制均由后端执行，前端提供同等或更早的提示。长度按用户输入字符计，前后端采用一致的计数规则；数据库字段长度不得小于允许的输入长度。所有必填文本在去除首尾空白后检查非空，密码除外。

必填性按对应操作判定。商品及内容草稿允许省略发布时必填的字段；已填写值仍须满足类型、长度、范围和唯一性要求。发布及已发布内容编辑时应完整校验，不能通过草稿规则绕过发布约束。

\begingroup
\smalltable
\begin{longtable}{>{\raggedright\arraybackslash}p{\dimexpr 0.200\linewidth-2\tabcolsep\relax}>{\raggedright\arraybackslash}p{\dimexpr 0.480\linewidth-2\tabcolsep\relax}>{\raggedright\arraybackslash}p{\dimexpr 0.320\linewidth-2\tabcolsep\relax}}
\caption{数据字典与输入约束}\label{tab:17}\\
\toprule
{\bfseries 字段或数据} & {\bfseries 本版本约束} & {\bfseries 错误提示示例} \\
\midrule
\endfirsthead
\multicolumn{3}{r}{\footnotesize 表~\thetable~（续）}\\
\toprule
{\bfseries 字段或数据} & {\bfseries 本版本约束} & {\bfseries 错误提示示例} \\
\midrule
\endhead
\midrule
\multicolumn{3}{r}{\footnotesize 续下页}\\
\endfoot
\bottomrule
\endlastfoot
邮箱 & 非空、合法邮箱结构、最多 254 字符；注册按规范化后唯一校验 & “请输入有效邮箱”／“该邮箱已注册”。 \\
\tabledivider
密码 & 8--64 字符，含字母和数字，确认密码一致 & “密码需为 8--64 个字符，并包含字母和数字”。 \\
\tabledivider
昵称／收件人／联系人 & 昵称 2--30；收件人和联系人 2--50 字符 & “请填写 2--50 个字符的收件人姓名”。 \\
\tabledivider
电话 & 去除允许的空格、连字符和括号后，允许开头一个 \code{+}，其余为 6--20 位数字 & “请输入有效联系电话，可包含国家区号”。 \\
\tabledivider
国家／地区、城市 & 结算和申请必填；采用受控选项或 2--100 字符文本 & “请填写收货国家／地区和城市”。 \\
\tabledivider
详细地址 & 5--200 字符；仅用于成人收货联系人 & “请填写完整的收货地址”。 \\
\tabledivider
商品名／SKU & 名称 2--100；SKU 3--40 个大写字母、数字或连字符，唯一且创建后固定 & “SKU 已存在，请使用其他编号”。 \\
\tabledivider
价格 & 大于 0，最多两位小数，单价不超过 999999.99 CNY & “价格必须是大于 0 的金额，最多保留两位小数”。 \\
\tabledivider
库存 & 整数，不小于 0；增加和减少为正整数，调整后不越界 & “本次调整将使库存不足，请减少调整数量”。 \\
\tabledivider
年龄 & 最小年龄 0--18 的整数；最大年龄可空，非空时为不小于最小年龄且不超过 18 的整数 & “年龄上限不能小于下限”；无上限用于“6 岁及以上”等商品。 \\
\tabledivider
零售数量 & 每行 1--99 的整数且不超过可售库存；购物车不超过 20 行 & “购买数量必须为 1--99 的整数”。 \\
\tabledivider
最小询价数量与询价数量 & 最小量 1--9999；询价数量不低于最小量且不超过 9999；每单不超过 20 行 & “该商品最少询价数量为 12 件”。 \\
\tabledivider
期望交付日期 & 可空；填写时不早于提交当天，按展示时区判断 & “期望日期不能早于今天”。 \\
\tabledivider
商品文案 & 简述 1--200 字符，详细说明最多 10000；玩法、安全提示、包装内容和规格说明各 1--2000 & “请补充商品安全提示”。 \\
\tabledivider
内容标题／正文 & 文章标题 2--100 字符、发布正文 1--20000；FAQ 问题 2--200、答案 1--5000；Banner 标题 1--100 & “发布前请填写标题和正文”。 \\
\tabledivider
主题／咨询正文／询价说明 & 主题 2--100；咨询正文 10--2000；询价交付说明、用途和备注均可空，各为 0--2000 & “内容最多允许 2000 个字符”。 \\
\tabledivider
回复／审核或调整原因 & 公开回复 1--2000；需填写的原因 2--500；内部备注最多 2000 & “请填写驳回原因，申请人将看到这段说明”。 \\
\tabledivider
网站和跳转链接 & 允许合法的站内路径或 HTTP／HTTPS 地址；最多 2048 字符，禁止脚本协议 & “请输入有效网页地址”。 \\
\tabledivider
发货信息 & 模拟物流名称 2--50；模拟运单号 3--50 个字母、数字或连字符 & “请填写模拟物流名称和运单号”。 \\
\tabledivider
上传文件 & 图片 $\leq$ 5 MiB；PDF $\leq$ 10 MiB；同时验证内容类型 & “仅支持 PDF，且文件不能超过 10 MiB”。 \\
\tabledivider
列表参数 & 页码为正整数；每页 1--50；排序和状态来自允许值 & “分页或排序参数无效”。 \\
\end{longtable}
\endgroup

可选字段未填写时按空值处理，不用“无”“未知”等占位文本冒充真实数据。未定义的接口输入字段应拒绝或忽略且无副作用，不能通过通用对象赋值更新角色、归属、价格快照或状态。

\Needspace{10\baselineskip}
\subsection{质量属性}\label{sec:35-质量属性}

\reqheading[3]{NFR-01}{NFR-01 业务一致性与持久化}

系统应持久化保存 P0 功能产生的数据。刷新页面、退出再登录或重启服务后，已成功提交的购物车、订单、申请、工单、询价及运营变更应保持可查询；不得仅保存在客户端内存中。业务写入失败时应保持原有一致状态，相关原子性和恢复要求按 BR-06 执行。

{\bfseries 业务理由}：保证业务处理结果在会话和服务生命周期之外持续有效。{\bfseries 验证方法}：TC-12、TC-17、TC-40。

\reqheading[3]{NFR-02}{NFR-02 安全与隐私}

系统应使用成熟的自适应密码摘要方案，不保存明文或可逆密码；数据库查询应参数化。受保护接口和文件应执行服务端授权，公开网络部署应使用 HTTPS；密钥、会话凭证和管理员初始密码不得暴露在前端或源代码仓库。

采用 Cookie 会话时，应设置 HttpOnly、适当的 SameSite 属性，并防护跨站写请求；采用令牌方案时，退出、停用和过期同样应阻止继续访问。输入转义、富文本清洗和文件验证按第 3.1.3 节、FR-29 执行。

登录限制按 FR-10 执行。注册接口每个来源地址最多允许 5 次请求／10 分钟；咨询、申请、询价写接口每个已登录用户合计最多允许 20 次请求／10 分钟。超限应返回等待提示，不新增业务记录；上线配置变更应同步记录并重新验证。

系统应遵循第 2.3.1 节的成人账户与数据采集范围，明确显示数据用途、必要确认及服务联系渠道。公开渠道信息须经独立发布，企业申请资料及内部备注不得自动公开。非公开数据最小化返回，关键操作记录不得包含密码或会话凭证。

{\bfseries 业务理由}：保护账户、个人信息及经销合作资料，降低未授权访问与恶意提交的影响。{\bfseries 验证方法}：TC-06、TC-08 至 TC-11、TC-26、TC-29、TC-32、TC-35。

\reqheading[3]{NFR-03}{NFR-03 响应式与可用性}

系统应在 390 px、768 px、1440 px 三种页面宽度下正确呈现首页、商品列表、详情、购物车、结算、申请表及后台订单页。主要内容和操作不得被遮挡，页面整体不得横向溢出；后台宽表允许在明确的局部区域滚动。

注册、登录、搜索、结算、申请和咨询表单应具有可见标签、可识别的键盘焦点与文本错误提示。核心交互应支持键盘操作，错误不得仅由颜色表达，图片应具有内容描述或恰当的装饰标记。

{\bfseries 业务理由}：保证消费者、采购人员和平台人员能够在常用设备上完成任务。{\bfseries 验证方法}：TC-01、TC-09、TC-38、TC-41。

\reqheading[3]{NFR-04}{NFR-04 性能与兼容性}

性能基线采用第 5.1 节数据量、单客户端和无外部服务依赖的验收环境。普通商品列表、详情和本人订单接口预热后各连续请求 20 次，每类至少 19 次在 1 秒内完成响应。首页、商品列表、详情及购物车分别加载 5 次，每页至少 4 次在 3 秒内显示主要数据并可操作。文件传输耗时另行记录。

测试记录应包括服务器和客户端配置、网络条件、数据库及浏览器版本、数据量、测量起止点和原始结果。至少在 Chrome 以及 Safari 或 Edge 中的一种浏览器分别完成关键流程验证，总计不少于两种浏览器。上述指标不等同于高并发吞吐量或商业服务可用率承诺；容量目标与运行环境由部署设计明确。

{\bfseries 业务理由}：为常用查询和交互提供可复现的性能及兼容性基线。{\bfseries 验证方法}：TC-38、TC-39。

\Needspace{10\baselineskip}
\subsection{实施与交付约束}\label{sec:36-实施与交付约束}

\reqheading[3]{NFR-05}{NFR-05 可部署性与可验证性}

系统应提供前端、业务服务和数据库的安装启动说明、依赖版本、环境变量示例、初始化方式、测试数据准备方式、管理员初始化方式及验证命令。具体命令应在选定技术方案并执行验证后写入运行文档；数据库和文件存储位置应可配置。

订单、授权、申请审核和库存逻辑应能够通过可重复执行的接口测试或脚本验证。测试库和测试账户应可独立初始化，重置操作不得影响业务运行数据；验收不得依赖手工修改数据库制造中间状态。

交付物应包括可运行的软件及源代码、数据库初始化与升级说明、接口说明、运行配置说明、测试程序和验证记录。版本变更应保持需求、实现与验证场景之间可追踪；技术方案与部署环境的变化不得绕过本规格规定的业务约束。

{\bfseries 业务理由}：支持可靠部署、问题复现及后续维护。{\bfseries 验证方法}：TC-40、配置与交付文档检查。
